\section[GIROSCOPÍA Y TRANSFORMACIONES CINEMÁTICAS]{CINEMÁTICA ROTATORIA, GIROSCOPÍA Y TRANSPORTE ENTRE MARCOS}
\label{prefgyro:section}

Esta sección establece una única cadena cinemática para la masa distribuida de
las palas, el cubo, las restricciones rígidas raíz--cubo y la transformación
multiblade. El punto de partida es una rotación nominal de velocidad angular
constante, impuesta alrededor de un eje conocido, y pequeñas perturbaciones
respecto del equilibrio elástico predeformado. La hipótesis de giro constante es
esencial: no se identifica el movimiento nominal con el de un sólido libre. Si
el eje prescrito no es principal, existe un par de mantenimiento que debe formar
parte de las acciones del punto operativo.

La separación entre movimiento nominal y perturbación sigue la construcción de
pequeñas oscilaciones alrededor de un movimiento estacionario descrita en la
Sección~4.4, ecuaciones (4.51)--(4.70), de \citet{meirovitch1997}. Las matrices
seccionales que se obtienen a continuación son, sin embargo, la especialización
particular adoptada aquí para una sección de viga con traslación, giro y centro
de masa excéntrico; no se atribuyen como una fórmula literal de esa referencia.

\subsection{Energía cinética exacta de una sección}
\label{prefgyro:exacta}

Sea $\bm R_0(t)\in SO(3)$ la orientación nominal del marco que gira con el
rotor y sea

\begin{equation}
 \bm W=[\bm\Omega]_\times,
 \qquad
 \dot{\bm R}_0=\bm R_0\bm W,
 \qquad
 \dot{\bm\Omega}=\bm 0 .
 \label{prefgyro:W}
\end{equation}

Las componentes de $\bm\Omega$ se expresan en el marco nominal. Para una
sección, $\bm a$ es el vector desde el eje de giro hasta el punto de referencia
de la línea de viga, $\bm r$ es el offset desde ese punto hasta el centro de
masa, $m$ es la masa por unidad de longitud y $\bm J$ es el tensor de inercia
centroidal por unidad de longitud. Se introducen la traslación perturbada
$\bm u$ y el vector de rotación perturbada $\bm\theta$ mediante

\begin{equation}
 \bm R(t)=\bm R_0(t)\bm E(\bm\theta),
 \qquad
 \bm E(\bm\theta)=\exp([\bm\theta]_\times).
 \label{prefgyro:Rexacta}
\end{equation}

La posición exacta del centro de masa, salvo una traslación común del origen,
es

\begin{equation}
 \bm x_G=\bm R_0
 \left(\bm a+\bm u+\bm E\bm r\right).
 \label{prefgyro:xG}
\end{equation}

Si $\bm J_\ell(\bm\theta)$ es el Jacobiano izquierdo de $SO(3)$,

\begin{equation}
 \bm J_\ell(\bm\theta)
 =\bm I+\frac{1-\cos\vartheta}{\vartheta^2}[\bm\theta]_\times
 +\frac{\vartheta-\sin\vartheta}{\vartheta^3}
 [\bm\theta]_\times^2,
 \qquad \vartheta=\lVert\bm\theta\rVert,
 \label{prefgyro:Jleft}
\end{equation}

entonces
$\dot{\bm E}\bm E^{\mathsf T}
 =[\bm J_\ell(\bm\theta)\dot{\bm\theta}]_\times$.
Al expresar las velocidades en el marco nominal, donde la rotación
$\bm R_0$ no altera las normas, se obtiene

\begin{align}
 \bm v_G={}&
 \bm W\left(\bm a+\bm u+\bm E\bm r\right)
 +\dot{\bm u}
 +[\bm J_\ell(\bm\theta)\dot{\bm\theta}]_\times
 \bm E\bm r,
 \label{prefgyro:vexacta}\\
 \bm\omega_G={}&
 \bm\Omega+\bm J_\ell(\bm\theta)\dot{\bm\theta}.
 \label{prefgyro:omegaexacta}
\end{align}

La energía cinética exacta por unidad de longitud es, por tanto,

\begin{equation}
 \mathcal T_{\mathrm{exact}}
 =\frac12m\bm v_G^{\mathsf T}\bm v_G
 +\frac12\bm\omega_G^{\mathsf T}
 \bm E\bm J\bm E^{\mathsf T}\bm\omega_G .
 \label{prefgyro:Texacta}
\end{equation}

Esta expresión retiene simultáneamente la velocidad orbital, la velocidad de la
línea de referencia, la excentricidad del centro de masa y la inercia rotacional
anisótropa. En particular, no es correcto añadir a posteriori un bloque de
momento angular a una masa traslacional sin verificar que ambos provienen de la
misma expansión de \eqref{prefgyro:Texacta}.

\subsection{Expansión de segundo orden y matrices seccionales completas}
\label{prefgyro:segundo-orden}

Para obtener las ecuaciones linealizadas se emplean

\begin{align}
 \bm E(\bm\theta)
 & =\bm I+[\bm\theta]_\times
 +\frac12[\bm\theta]_\times^2
 +\mathcal O(\lVert\bm\theta\rVert^3),
 \label{prefgyro:exp2}\\
 \bm J_\ell(\bm\theta)\dot{\bm\theta}
 & =\dot{\bm\theta}
 +\frac12\bm\theta\times\dot{\bm\theta}
 +\mathcal O(3).
 \label{prefgyro:Jleft2}
\end{align}

Resulta útil definir, para dos vectores cualesquiera $\bm a$ y $\bm b$, la
matriz simétrica

\begin{equation}
 \mathcal H(\bm a,\bm b)
 =\frac12\left(\bm a\bm b^{\mathsf T}
 +\bm b\bm a^{\mathsf T}\right)
 -(\bm a^{\mathsf T}\bm b)\bm I,
 \label{prefgyro:Hab}
\end{equation}

que satisface la identidad

\begin{equation}
 \bm a^{\mathsf T}[\bm\theta]_\times^2\bm b
 =\bm\theta^{\mathsf T}\mathcal H(\bm a,\bm b)\bm\theta.
 \label{prefgyro:Hidentity}
\end{equation}

Se agrupan las seis perturbaciones seccionales en
$\bm\eta=[\bm u^{\mathsf T},\bm\theta^{\mathsf T}]^{\mathsf T}$ y se definen

\begin{equation}
 \bm B_G=\begin{bmatrix}\bm I&-[\bm r]_\times\end{bmatrix},
 \quad
 \bm b=\bm a+\bm r,
 \quad
 \bm p=m\bm W\bm b,
 \quad
 \bm h=\bm J\bm\Omega .
 \label{prefgyro:definiciones}
\end{equation}

Aquí $\bm p$ y $\bm h$ son, respectivamente, el momento lineal orbital del
centro de masa y el momento angular centroidal del estado nominal. La parte
cuadrática de la energía puede escribirse como

\begin{equation}
 \mathcal T^{(2)}
 =\frac12\dot{\bm\eta}^{\mathsf T}\bm M_6\dot{\bm\eta}
 +\dot{\bm\eta}^{\mathsf T}\bm A_6\bm\eta
 +\frac12\bm\eta^{\mathsf T}\bm H_\Omega\bm\eta .
 \label{prefgyro:T2}
\end{equation}

La masa seccional completa es

\begin{equation}
 \bm M_6
 =m\bm B_G^{\mathsf T}\bm B_G
 +\begin{bmatrix}\bm 0&\bm 0\\\bm 0&\bm J\end{bmatrix}
 =\begin{bmatrix}
 m\bm I&-m[\bm r]_\times\\
 m[\bm r]_\times&\bm J-m[\bm r]_\times^2
 \end{bmatrix}.
 \label{prefgyro:M6}
\end{equation}

La matriz completa de acoplamiento velocidad--posición es

\begin{equation}
 \bm A_6
 =\begin{bmatrix}
 m\bm W&-m\bm W[\bm r]_\times\\
 m[\bm r]_\times\bm W&
 -m[\bm r]_\times\bm W[\bm r]_\times
 +\mathcal H(\bm p,\bm r)+\bm J\bm W-\tfrac12[\bm h]_\times
 \end{bmatrix}.
 \label{prefgyro:A6}
\end{equation}

El término $\mathcal H(\bm p,\bm r)$ es simétrico y, por consiguiente, no
aparece en la matriz giroscópica. Se conserva para presentar la expansión
completa: con $\bm p,\bm r$ constantes, su contribución es la derivada total
$\frac{\dd}{\dd t}[\tfrac12\bm\theta^{\mathsf T}\mathcal H(\bm p,\bm r)\bm\theta]$
y puede omitirse sin modificar las ecuaciones. La contribución de
velocidad en las ecuaciones de Lagrange es

\begin{equation}
 \boxed{\bm G_6=\bm A_6-\bm A_6^{\mathsf T}}
 =\begin{bmatrix}
 2m\bm W&-2m\bm W[\bm r]_\times\\
 2m[\bm r]_\times\bm W&
 -2m[\bm r]_\times\bm W[\bm r]_\times
 +\bm J\bm W+\bm W\bm J-[\bm h]_\times
 \end{bmatrix}.
 \label{prefgyro:G6}
\end{equation}

Por construcción, $\bm G_6^{\mathsf T}=-\bm G_6$. El bloque rotacional
$\bm J\bm W+\bm W\bm J-[\bm J\bm\Omega]_\times$ es necesario incluso cuando
el centro de masa está sobre la línea de referencia. El bloque reducido
$[\bm J\bm\Omega]_\times$ solo representa una parte asociada al momento
angular nominal: no contiene los términos $\bm J\bm W+\bm W\bm J$ y, por
tanto, no es la expansión completa de \eqref{prefgyro:Texacta}.

La anisotropía muestra la diferencia con claridad. Si $\bm\Omega=\Omega\bm e_1$
y $\bm J=\operatorname{diag}(J_1,J_2,J_3)$ en ejes principales, el acoplamiento
en el plano $2$--$3$ queda multiplicado por

\begin{equation}
 (J_2+J_3-J_1)\Omega,
 \label{prefgyro:principal}
\end{equation}

y no por $J_1\Omega$. Para una inercia no diagonal deben conservarse todos los
productos de inercia de \eqref{prefgyro:G6}.

La Hessiana puramente centrífuga que aparece en \eqref{prefgyro:T2} es

\begin{equation}
 \bm H_\Omega
 =\begin{bmatrix}
 m\bm W^{\mathsf T}\bm W&
 -m\bm W^{\mathsf T}\bm W[\bm r]_\times\\
 m[\bm r]_\times\bm W^{\mathsf T}\bm W&
 -m[\bm r]_\times\bm W^{\mathsf T}\bm W[\bm r]_\times
 +\mathcal H(\bm W^{\mathsf T}\bm p,\bm r)
 +\mathcal H(\bm\Omega,\bm h)
 -\bm W\bm J\bm W
 \end{bmatrix}.
 \label{prefgyro:HOmega}
\end{equation}

La fuerza generalizada constante inducida por el giro se obtiene de la parte
lineal de la energía:

\begin{equation}
 \bm f_{\Omega,0}
 =\bm B_G^{\mathsf T}m\bm W^{\mathsf T}\bm W\bm b
 +\begin{bmatrix}\bm 0\\ \bm h\times\bm\Omega\end{bmatrix}.
 \label{prefgyro:fcentrif}
\end{equation}

Con la convención de almacenar las fuerzas centrífugas en el lado derecho,

\begin{equation}
 \bm f_\Omega(\bm\eta)
 =\bm f_{\Omega,0}+\bm H_\Omega\bm\eta,
 \qquad
 \bm K_{\Omega,c}=-\bm H_\Omega
 \label{prefgyro:KOmega}
\end{equation}

es la contribución correspondiente en el lado izquierdo. Debe sumarse a la
rigidez elástica y a la rigidez geométrica del equilibrio prestresado, no a una
configuración descargada elegida de manera independiente.

Cuando $\bm h$ no es paralelo a $\bm\Omega$, la rotación uniforme prescrita
requiere el par

\begin{equation}
 \bm\tau_{\mathrm{hold}}=\bm\Omega\times\bm h.
 \label{prefgyro:holding}
\end{equation}

Este par se opone al bloque $\bm h\times\bm\Omega$ de
\eqref{prefgyro:fcentrif} y compensa el residuo nominal de Euler. Si no se
aplica, el eje de giro no permanece constante y las hipótesis de
\eqref{prefgyro:W} dejan de ser
válidas. Por ello las ecuaciones anteriores describen un rotor mantenido en su
punto operativo; no constituyen una afirmación sobre la rotación libre de un
sólido anisótropo.

La fuerza centrífuga y el par de mantenimiento participan en el equilibrio de
la turbina completa. Después del ensamblaje, dicho equilibrio tiene la forma

\begin{equation}
 \bm T_0^{\mathsf T}\left[
 \bm f_{\mathrm{int}}(\bm q_0)
 -\bm f_{\mathrm{ext}}^{\mathrm{nc}}(\bm q_0)
 -\bm f_{\Omega}(\bm q_0)\right]=\bm 0,
 \label{prefgyro:equilibrio-nominal}
\end{equation}

donde $\bm f_{\Omega}(\bm q_0)=\bm f_{\Omega,0}+\bm H_\Omega\bm q_0$,
$\bm f_{\mathrm{ext}}^{\mathrm{nc}}$ agrupa las acciones no centrífugas y
$\bm f_{\mathrm{int}}$ incluye las fuerzas elásticas del tren y de la base.
Las reacciones de restricción se eliminan por proyección. No corresponde imponer
$\bm f_{\Omega,0}=\bm 0$ como condición de un rotor libre. Además, la tangente
del par depende de cómo se mantenga la carga durante la perturbación. Con la
convención de residual en el lado izquierdo y coordenadas exponenciales, dos
casos representativos son

\begin{equation}
 \bm K_{\tau}^{\mathrm{cuerpo}}
 =+\frac12[\bm\tau_{\mathrm{hold}}]_\times,
 \qquad
 \bm K_{\tau}^{\mathrm{nominal}}
 =-\frac12[\bm\tau_{\mathrm{hold}}]_\times.
 \label{prefgyro:holding-tangent}
\end{equation}

El primer signo corresponde a componentes constantes en el cuerpo perturbado y
el segundo a componentes constantes en el marco nominal. Si el par proviene de
una ley elástica o de control del tren, debe utilizarse la derivada de esa ley,
en lugar de adoptar uno de estos dos tangentes como regla universal.

Para un elemento de viga, las matrices seccionales se proyectan mediante la
misma interpolación cinemática $\bm N(x)$ usada por la masa consistente:

\begin{equation}
 \bm M_e=\int_0^L\bm N^{\mathsf T}\bm M_6\bm N\,\dd x,
 \quad
 \bm G_{c,e}=\int_0^L\bm N^{\mathsf T}\bm G_6\bm N\,\dd x,
 \quad
 \bm K_{\Omega,c,e}=-\int_0^L
 \bm N^{\mathsf T}\bm H_\Omega\bm N\,\dd x.
 \label{prefgyro:elemento}
\end{equation}

La dependencia de $\bm a(x)$ debe conservarse al evaluar
$\bm f_{\Omega,0}$ y $\bm H_\Omega$. Esta integración común evita combinar una
masa consistente con una giroscopía o una fuerza centrífuga de interpolación
distinta.

\subsection{Transporte de las perturbaciones al marco inercial}
\label{prefgyro:transporte}

Las matrices anteriores están expresadas en componentes corrotantes. Para
acoplar el rotor con la góndola y la torre se transportan todas las
perturbaciones nodales del rotor a ejes inerciales antes de imponer las
restricciones. Sea $\bm q_c=\bm P^{\mathsf T}(t)\bm q_g$ y, en un elemento de
dos nodos,

\begin{equation}
 \bm{\mathcal W}_{12}
 =\operatorname{blkdiag}(\bm W,\bm W,\bm W,\bm W).
 \label{prefgyro:W12}
\end{equation}

Antes de aplicar estas relaciones, las matrices y $\bm W$ se expresan en una misma base de componentes de ensamblaje. En el azimut instantáneo puede elegirse $\bm P=\bm I$ sin perder las derivadas
temporales. Para giro constante,

\begin{equation}
 \dot{\bm q}_c=\dot{\bm q}_g-\bm{\mathcal W}_{12}\bm q_g,
 \qquad
 \ddot{\bm q}_c=\ddot{\bm q}_g
 -2\bm{\mathcal W}_{12}\dot{\bm q}_g
 +\bm{\mathcal W}_{12}^2\bm q_g.
 \label{prefgyro:qtransport}
\end{equation}

Al sustituir \eqref{prefgyro:qtransport} en la ecuación corrotante se obtiene

\begin{align}
 \bm D_g^{\mathrm{in}}&=\bm G_c-2\bm M_e\bm{\mathcal W}_{12},
 \label{prefgyro:Dg-in}\\
 \dot{\bm M}_g&=
 \bm{\mathcal W}_{12}\bm M_e
 -\bm M_e\bm{\mathcal W}_{12},
 \label{prefgyro:Mdot}\\
 \bm G_g&=\bm D_g^{\mathrm{in}}-\dot{\bm M}_g
 =\bm G_c-\bm M_e\bm{\mathcal W}_{12}
 -\bm{\mathcal W}_{12}\bm M_e,
 \label{prefgyro:Gg}\\
 \bm K_{\mathrm{trans}}&=
 -\bm G_c\bm{\mathcal W}_{12}
 +\bm M_e\bm{\mathcal W}_{12}^2.
 \label{prefgyro:Ktransport}
\end{align}

La matriz $\dot{\bm M}_g$ es simétrica y $\bm G_g$ es antisimétrica. En
cambio, el operador total de velocidad

\begin{equation}
 \bm D_g=\bm C_g+\dot{\bm M}_g+\bm G_g
 \label{prefgyro:Dtotal}
\end{equation}

no tiene por qué ser antisimétrico. Exigir antisimetría a $\bm D_g$ eliminaría
precisamente la derivada de la masa producida por una sección anisótropa o por
una interpolación cuya base gira. La rigidez del punto operativo en ejes
inerciales es

\begin{equation}
 \bm K_g^{\mathrm{op}}
 =\bm K_e+\bm K_{\mathrm{prestress}}
 +\bm K_{\Omega,c}+\bm K_{\mathrm{trans}}.
 \label{prefgyro:Koperativa}
\end{equation}

$\bm K_{\mathrm{trans}}$ es un término de transporte de coordenadas y no una
segunda copia de la rigidez centrífuga. Para aceleración angular no nula aparece
además $-\bm M_e\dot{\bm{\mathcal W}}_{12}$ en
\eqref{prefgyro:Ktransport}; dicha extensión queda fuera del punto operativo
de velocidad constante considerado aquí.

Para un cuerpo rígido con centro de masa en el maestro, las mismas identidades
dan, en el bloque rotacional inercial,

\begin{equation}
 \dot{\bm J}_g=\bm W\bm J_g-\bm J_g\bm W,
 \qquad
 \bm D_{\theta,g}=\dot{\bm J}_g-[\bm J_g\bm\Omega]_\times,
 \qquad
 \bm G_{\theta,g}=-[\bm J_g\bm\Omega]_\times.
 \label{prefgyro:rigidbody-inertial}
\end{equation}

Así, una masa cuyo centro coincide con el maestro no recibe un término de
Coriolis traslacional en coordenadas inerciales. El efecto angular del spin
prescrito sí permanece. Para un tensor anisótropo, $\dot{\bm J}_g$ debe
conservarse y el coeficiente completo de velocidad no es puramente
antisimétrico.

\subsection{Restricción raíz--cubo dependiente del tiempo}
\label{prefgyro:mpc}

Las raíces de pala se vinculan al maestro del cubo después del transporte a
componentes inerciales. Si

\begin{equation}
 \bm r_s(t)=\bm R_\Omega(\psi)\bm r_{s0},
 \qquad
 \dot{\bm r}_s=\bm W\bm r_s,
 \qquad
 \ddot{\bm r}_s=\bm W^2\bm r_s,
 \label{prefgyro:rslave}
\end{equation}

la compatibilidad infinitesimal entre maestro $m$ y esclavo $s$ es

\begin{equation}
 \begin{bmatrix}\bm u_s\\\bm\theta_s\end{bmatrix}
 =\bm B_s(t)
 \begin{bmatrix}\bm u_m\\\bm\theta_m\end{bmatrix},
 \qquad
 \bm B_s(t)=
 \begin{bmatrix}
 \bm I&-[\bm r_s(t)]_\times\\
 \bm 0&\bm I
 \end{bmatrix}.
 \label{prefgyro:Bslave}
\end{equation}

Por tanto,

\begin{equation}
 \dot{\bm B}_s=
 \begin{bmatrix}\bm 0&-[\dot{\bm r}_s]_\times\\\bm 0&\bm 0\end{bmatrix},
 \qquad
 \ddot{\bm B}_s=
 \begin{bmatrix}\bm 0&-[\ddot{\bm r}_s]_\times\\\bm 0&\bm 0\end{bmatrix}.
 \label{prefgyro:Bslave-derivadas}
\end{equation}

En las ecuaciones siguientes, el subíndice $f$ denota los operadores completos, anteriores a la reducción de restricciones. Al ensamblar todos los bloques se obtiene
$\bm q_f=\bm T_c(t)\bm q_r$. La reducción dinámica exacta es

\begin{align}
 \bm M_r&=\bm T_c^{\mathsf T}\bm M_f\bm T_c,
 \label{prefgyro:Mred}\\
 \bm D_r&=\bm T_c^{\mathsf T}
 \left(\bm D_f\bm T_c+2\bm M_f\dot{\bm T}_c\right),
 \label{prefgyro:Dred}\\
 \bm K_r&=\bm T_c^{\mathsf T}
 \left(\bm K_f\bm T_c+\bm D_f\dot{\bm T}_c
 +\bm M_f\ddot{\bm T}_c\right).
 \label{prefgyro:Kred}
\end{align}

Una congruencia estática $\bm T_c^{\mathsf T}(\cdot)\bm T_c$ solo es suficiente
cuando $\dot{\bm T}_c=\ddot{\bm T}_c=\bm 0$. La derivada de la masa reducida,

\begin{equation}
 \dot{\bm M}_r
 =\dot{\bm T}_c^{\mathsf T}\bm M_f\bm T_c
 +\bm T_c^{\mathsf T}\dot{\bm M}_f\bm T_c
 +\bm T_c^{\mathsf T}\bm M_f\dot{\bm T}_c,
 \label{prefgyro:Mred-dot}
\end{equation}

permite separar nuevamente la parte simétrica dependiente del tiempo y la
parte giroscópica antisimétrica. Esta construcción mantiene la compatibilidad
en la interfaz sin asignar marcos distintos a la raíz y al cubo.

\subsection{Marco de pala temporal y transformación multiblade}
\label{prefgyro:mbc}

Después de la reducción de las raíces, los grados de libertad flexibles de las
tres palas se expresan en marcos homólogos. Si $\bm Q_{b0}$ contiene como
columnas los ejes locales de la pala $b$ en el azimut de referencia, el cambio
de base local--inercial es

\begin{equation}
 \bm Q_b(\psi)=\bm R_\Omega(\psi)\bm Q_{b0},
 \qquad
 \dot{\bm Q}_b=\bm W\bm Q_b,
 \qquad
 \ddot{\bm Q}_b=\bm W^2\bm Q_b.
 \label{prefgyro:Qblade}
\end{equation}

La matriz $\bm H(\psi)$ se construye colocando
$\operatorname{blkdiag}(\bm Q_b,\bm Q_b)$ en las traslaciones y rotaciones de
cada estación no raíz. En consecuencia,

\begin{equation}
 \dot{\bm H}=\bm{\mathcal W}_H\bm H,
 \qquad
 \ddot{\bm H}=\bm{\mathcal W}_H^2\bm H,
 \label{prefgyro:Htemporal}
\end{equation}

donde $\bm{\mathcal W}_H$ contiene $\bm W$ en los bloques rotantes y ceros en
los grados de libertad estacionarios. Tomar $\dot{\bm H}=\bm 0$ mezclaría una
base local corrotante con perturbaciones declaradas inerciales.

Sea $\bm B_{\mathrm{MBC}}(\psi)$ la transformación armónica colectiva--coseno--
seno para tres palas, con azimuts
$\psi_b=\psi+2\pi(b-1)/3$ \citep{bir2008}. Para un grado de libertad homólogo
por pala, su convención directa y su inversa son

\begin{equation}
 \bm B_{\mathrm{MBC}}=
 \begin{bmatrix}
 1&\cos\psi_1&\sin\psi_1\\
 1&\cos\psi_2&\sin\psi_2\\
 1&\cos\psi_3&\sin\psi_3
 \end{bmatrix},
 \qquad
 \bm B_{\mathrm{MBC}}^{-1}=\frac13
 \begin{bmatrix}
 1&1&1\\
 2\cos\psi_1&2\cos\psi_2&2\cos\psi_3\\
 2\sin\psi_1&2\sin\psi_2&2\sin\psi_3
 \end{bmatrix}.
 \label{prefgyro:Bmbc}
\end{equation}

Con $c_b=\cos\psi_b$, $s_b=\sin\psi_b$ y
$\omega_s=\dot\psi$ constante, las filas de sus derivadas son

\begin{equation}
 [\bm B_{\mathrm{MBC}}]_{b,:}=[1,c_b,s_b],\qquad
 [\dot{\bm B}_{\mathrm{MBC}}]_{b,:}
 =\omega_s[0,-s_b,c_b],\qquad
 [\ddot{\bm B}_{\mathrm{MBC}}]_{b,:}
 =-\omega_s^2[0,c_b,s_b].
 \label{prefgyro:Bmbc-derivadas}
\end{equation}

Para varios grados de libertad homólogos se emplea el producto de Kronecker
con la identidad correspondiente y se conserva la identidad en los grados de
libertad estacionarios. La transformación completa es

\begin{equation}
 \bm P_{\mathrm{MBC}}(\psi)
 =\bm H(\psi)\bm B_{\mathrm{MBC}}(\psi),
 \label{prefgyro:Pmbc}
\end{equation}

y sus derivadas deben conservar ambos orígenes temporales:

\begin{align}
 \dot{\bm P}_{\mathrm{MBC}}
 &=\dot{\bm H}\bm B_{\mathrm{MBC}}
 +\bm H\dot{\bm B}_{\mathrm{MBC}},
 \label{prefgyro:Pdot}\\
 \ddot{\bm P}_{\mathrm{MBC}}
 &=\ddot{\bm H}\bm B_{\mathrm{MBC}}
 +2\dot{\bm H}\dot{\bm B}_{\mathrm{MBC}}
 +\bm H\ddot{\bm B}_{\mathrm{MBC}}.
 \label{prefgyro:Pddot}
\end{align}

Finalmente, para $\bm q_r=\bm P_{\mathrm{MBC}}\bm z$, los operadores en
coordenadas multiblade son

\begin{align}
 \overline{\bm M}
 &=\bm P_{\mathrm{MBC}}^{\mathsf T}
 \bm M_r\bm P_{\mathrm{MBC}},
 \label{prefgyro:Mmbc}\\
 \overline{\bm D}
 &=\bm P_{\mathrm{MBC}}^{\mathsf T}
 \left(\bm D_r\bm P_{\mathrm{MBC}}
 +2\bm M_r\dot{\bm P}_{\mathrm{MBC}}\right),
 \label{prefgyro:Dmbc}\\
 \overline{\bm K}
 &=\bm P_{\mathrm{MBC}}^{\mathsf T}
 \left(\bm K_r\bm P_{\mathrm{MBC}}
 +\bm D_r\dot{\bm P}_{\mathrm{MBC}}
 +\bm M_r\ddot{\bm P}_{\mathrm{MBC}}\right).
 \label{prefgyro:Kmbc}
\end{align}

La coexistencia de $\bm G_g$ con las derivadas MBC no constituye por sí sola un
doble conteo. $\bm G_g$ procede de la energía y del transporte físico de las
perturbaciones, mientras que $\dot{\bm P}_{\mathrm{MBC}}$ y
$\ddot{\bm P}_{\mathrm{MBC}}$ aparecen al reexpresar esas mismas ecuaciones en
coordenadas armónicas. La consistencia exige conservar toda la cadena en el
orden presentado: expansión energética, transporte inercial, restricción
temporal raíz--cubo y, por último, transformación multiblade.
